% \documentclass[../main.tex]{subfiles}
% \begin{document}

\textbf{Notation.}
We denote by $\lambda$ the security parameter.
A function $f$ is \emph{negligible} in $\lambda$, denoted by $\mathrm{negl}(\lambda)$, if
$f(\lambda) < 1/\mathrm{poly}(\lambda)$ for every polynomial $\mathrm{poly}(\lambda)$ as $\lambda \to \infty$.
Computationally and statistically indistinguishable games are denoted by $\approx_c$ and $\approx_s$, respectively.
We denote the set $\{1,\dots,n\}$ by $[n]$, and the sequence $(x_{a_1}, \dots, x_{a_k})$ with
$a_1 < \cdots < a_k$ by $(x_{a_i})_{a_i \in \{a_1,\dots,a_k\}}$.
A positive semidefinite matrix $\mathbf{\Sigma}$ can be written as
$\mathbf{\Sigma} = \mathbf{S}\mathbf{S}^t$.
Therefore, we use $\sqrt{\mathbf{\Sigma}}$ to denote any fixed matrix $\mathbf{S}$ satisfying the above relation.

\section{Fundamental Lattice Theory}

\begin{definition}
Let $\mathbf{A} = [\mathbf{a}_1,\mathbf{a}_2,\dots,\mathbf{a}_m] \in \mathbb{R}^{n\times m}$ be a matrix whose columns are linearly independent.
The lattice generated by these vectors is defined as
\[
\Lambda(\mathbf{A})
=
\left\{
\sum_{i=1}^{m} z_i \mathbf{a}_i \,\middle|\, z_i \in \mathbb{Z}\ \forall i \in [m]
\right\}.
\]
The set $\{\mathbf{a}_1,\mathbf{a}_2,\dots,\mathbf{a}_m\}$ is called a basis of the lattice.
\end{definition}

\begin{definition}
Let $n,m,q$ be positive integers and let $\mathbf{A}\in\mathbb{Z}_q^{n\times m}$.
We define the following $m$-dimensional $q$-ary lattices associated with $\mathbf{A}$:
\[
\setlength{\jot}{2pt}
\begin{aligned}
\Lambda^\perp(\mathbf{A})
&=
\left\{
\mathbf{z}\in\mathbb{Z}^m : \mathbf{A}\mathbf{z}\equiv \mathbf{0}\pmod q
\right\},\\
\Lambda(\mathbf{A}^t)
&=
\left\{
\mathbf{v}\in\mathbb{Z}^m : \exists\,\mathbf{s}\in\mathbb{Z}_q^n \text{ such that } \mathbf{v}\equiv \mathbf{A}^t\mathbf{s}\pmod q
\right\},\\
\Lambda^{\mathbf{u}}(\mathbf{A})
&=
\left\{
\mathbf{v}\in\mathbb{Z}^m : \mathbf{A}\mathbf{v}\equiv \mathbf{u}\pmod q
\right\}
=
\mathbf{x} + \Lambda^\perp(\mathbf{A}),
\end{aligned}
\]
for $\mathbf{u} \in \mathbb{Z}_q^n$, where $\mathbf{x}\in\mathbb{Z}^m$ is any solution satisfying
$\mathbf{A}\mathbf{x}\equiv \mathbf{u}\pmod q$.
\end{definition}

Dual lattices play a central role in the analysis of Gaussian meansure, which is essential for the security proof of every cryptographic scheme


\begin{definition}
For a lattice $\Lambda$, we define the \emph{dual lattice}
$\Lambda^* = \{\mathbf{w} \in \mathbb{R}^n : \forall \mathbf{x} \in \Lambda,\ \mathbf{w}^{\top}\mathbf{x} \in \mathbb{Z}\}$.
\end{definition}
\noindent For all the matrix $\mathbf{A} \in \mathbb{Z}_q^{n \times m}$, it holds that $\Lambda(\mathbf{A}) = q \cdot \Lambda^\perp(\mathbf{A})^*$. Indeed, if we take $\mathbf{x} \in \Lambda(\mathbf{A})$, then for any $\mathbf{y} \in \Lambda^\perp(\mathbf{A})$, we have $\big \langle \dfrac{\mathbf{x}^\top}{q},\mathbf{y}\big\rangle= \big \langle \dfrac{\mathbf{z}^\top\mathbf{A}}{q},\mathbf{y}\big\rangle = \big \langle \dfrac{\mathbf{z}^\top}{q},\mathbf{A}\mathbf{y}\big\rangle=\big \langle \dfrac{\mathbf{z}^\top}{q},\mathbf{k}q\big\rangle =\big \langle \mathbf{z}^\top,\mathbf{k}\big\rangle \in \mathbb{Z}$ where $\mathbf{z}$ and $\mathbf{k}$ all come from $\mathbb{Z}^n$. And since this also holds for every $\mathbf{x} \in \Lambda(\mathbf{A})$, we have the desired result. 

For a lattice $\Lambda$, the smoothing parameter $\eta_{\varepsilon}(\Lambda)$ denotes the smallest Gaussian parameter $s$ such that the Gaussian mass over the nonzero dual lattice vectors becomes negligible, i.e.,
\[
\rho_{1/s}(\Lambda^* \setminus \{\mathbf{0}\}) \le \varepsilon.
\]
This parameter is a standard tool in lattice-based cryptography for analyzing discrete Gaussian sampling ~\cite{J:MicReg07}.

\begin{lemma}[From~\cite{C:GenPeiVai08}]
\label{lem:1}
For any $n$-dimensional lattice $\Lambda$ and real $\epsilon > 0$ and $\mathbf{S} \subset \Lambda$ is a basis of it, we have
\[
\eta_{\epsilon}(\Lambda)
\leq
\|\widetilde{S}\|
\cdot
\sqrt{\log(2n(1+1/\epsilon))/\pi}.
\]
Then, for any $\omega(\sqrt{\log n})$ function, there is a negligible
$\epsilon(n)$ for which
$
\eta_{\epsilon}(\Lambda)
\leq
\|\widetilde{S}\|\cdot \omega(\sqrt{\log n}).
$
\end{lemma}

\begin{lemma}[Gaussian Tail Bound~\cite{C:MicPei12}]
\label{lem:2}
Let $\mathbf{B} \in \mathbb{R}^{n \times n}$ be a full-rank basis and let
$\Lambda$ be the $n$-dimensional lattice generated by $\mathbf{B}$.
Suppose $s \ge \eta_{\varepsilon}(\Lambda)$ for some
$\varepsilon \in (0,1)$. Then, for all $\mathbf{c} \in \mathbb{R}^n$,
\[
\Pr\left[
    \|\mathbf{v}\| > s\sqrt{n}
    :
    \mathbf{v} \leftarrow \mathcal{D}_{\Lambda+\mathbf{c},s}
\right]
\le
2^{-n}\cdot \frac{1+\varepsilon}{1-\varepsilon}.
\]
\end{lemma}


\begin{definition}[Tensor Product / Kronecker Product]
Let $\mathbf{A} = (a_{i,j}) \in \mathbb{Z}^{n \times m}$ and $\mathbf{B} \in \mathbb{Z}^{\ell \times \lambda}$. The tensor product, also called the Kronecker product, of $\mathbf{A}$ and $\mathbf{B}$ is defined as
\[
\mathbf{A} \otimes \mathbf{B}
=
\begin{bmatrix}
a_{1,1}\mathbf{B} & \cdots & a_{1,m}\mathbf{B} \\
\vdots & \ddots & \vdots \\
a_{n,1}\mathbf{B} & \cdots & a_{n,m}\mathbf{B}
\end{bmatrix}
\in \mathbb{Z}^{n\ell \times m\lambda}.
\]
\end{definition}

The Kronecker product satisfies the following mixed-product property:
\[
(\mathbf{A} \otimes \mathbf{B})(\mathbf{C} \otimes \mathbf{D})
=
(\mathbf{A}\mathbf{C}) \otimes (\mathbf{B}\mathbf{D}),
\]
whenever the matrix products $\mathbf{A}\mathbf{C}$ and $\mathbf{B}\mathbf{D}$ are well-defined.

\begin{definition}
The $n$-dimensional Gaussian function is defined as
\[
\rho_{\mathbf{B},\mathbf{c}}(\mathbf{x})
=
\rho_{\sqrt{\mathbf{\Sigma}},\mathbf{c}}(\mathbf{x})
=
\exp\!\left(
-\pi (\mathbf{x}-\mathbf{c})^t \mathbf{\Sigma}^{-1} (\mathbf{x}-\mathbf{c})
\right),
\]
for invertible $\mathbf{B}\in\mathbb{R}^{n\times n}$, $\mathbf{c}\in\mathbb{R}^n$, and
$\mathbf{\Sigma}:=\mathbf{B}\mathbf{B}^t$.
When $\mathbf{\Sigma}=\sigma^2\mathbf{I}$ and $\mathbf{c}=\mathbf{0}$, we write $\rho_\sigma$.
\end{definition}

\begin{definition}
The discrete Gaussian distribution over $\Lambda$ with parameter $\sqrt{\mathbf{\Sigma}}$ and center $\mathbf{c}$ is defined as
\[
\mathcal{D}_{\Lambda,\mathbf{\Sigma},\mathbf{c}}(\mathbf{x})
=
\frac{\rho_{\sqrt{\mathbf{\Sigma}},\mathbf{c}}(\mathbf{x})}
{\sum_{\mathbf{y}\in\Lambda}\rho_{\sqrt{\mathbf{\Sigma}},\mathbf{c}}(\mathbf{y})}
=
\frac{\rho_{\sqrt{\mathbf{\Sigma}},\mathbf{c}}(\mathbf{x})}
{\rho_{\sqrt{\mathbf{\Sigma}},\mathbf{c}}(\Lambda)}.
\]
\end{definition}

\begin{definition}
    Let $X, Y$ are two discrete random variables over countable set $A$. The statistical distance between them is defined as 
    $$
    \triangle[X,Y] = \dfrac{1}{2}\sum_{a \in A}|\Pr\{X=a\}-\Pr\{X=b\}|
    $$
\end{definition}



\begin{lemma}[Leftover Hash Lemma~\cite{SIComp:HILL99}]
\label{lem:leftoverHash}
Let $n,m,q$ be lattice parameters, and suppose $m \ge 2n \log q$ and $q > 4$. Then, the statistical distance between the following distributions is at most $2^{-n}$:
\[
\left\{(\mathbf A,\mathbf A\mathbf x) : \mathbf A \xleftarrow{\$} \mathbb Z_q^{n\times m},\ \mathbf x \xleftarrow{\$} \{0,1\}^m \right\}
\quad \text{and} \quad
\left\{(\mathbf A,\mathbf y) : \mathbf A \xleftarrow{\$} \mathbb Z_q^{n\times m},\ \mathbf y \xleftarrow{\$} \mathbb Z_q^n \right\}.
\]

\end{lemma}
\begin{proof}
Let us define
$S_{\mathbf A}(\mathbf y)=\{\mathbf x \in \{0,1\}^m : \mathbf A\mathbf x=\mathbf y\}$
for a fixed matrix $\mathbf A$. Hence, the statistical distance between the given distributions is
$$
\begin{aligned}
\Delta
=
\mathbb E_{\mathbf A}
\left[
\sum_{\mathbf y \in \mathbb Z_q^n}
\left|
\dfrac{|S_{\mathbf A}(\mathbf y)|}{2^m}
-
\dfrac{1}{q^n}
\right|
\right]
&\leq
\mathbb E_{\mathbf A}
\left[
\sqrt{
q^n
\sum_{\mathbf y \in \mathbb Z_q^n}
\left(
\dfrac{|S_{\mathbf A}(\mathbf y)|}{2^m}
-
\dfrac{1}{q^n}
\right)^2
}
\right]
\\
&=
\mathbb E_{\mathbf A}
\left[
\sqrt{
q^n
\left(
\dfrac{1}{2^{2m}}
\sum_{\mathbf y \in \mathbb Z_q^n}
|S_{\mathbf A}(\mathbf y)|^2
-
\dfrac{1}{q^n}
\right)
}
\right],
\end{aligned}
$$
where the last equality follows from $\sum_{\mathbf y \in \mathbb Z_q^n}
|S_{\mathbf A}(\mathbf y)|
=
2^m.$

Assume we have $\mathbf x,\mathbf x' \in \{0,1\}^m$ and let
$\mathbf z = \mathbf x-\mathbf x' \in \{-1,0,1\}^m$.
We know that
$$
\Pr_{\mathbf A}
\left[
\mathbf A\mathbf z = \mathbf 0
:
\mathbf z \in \{-1,0,1\}^m,\ \mathbf z \neq \mathbf 0
\right]
=
\dfrac{1}{q^n}.
$$
Indeed, since $\mathbf z \neq \mathbf 0$, there exists some index $j$ such that $z_j=\pm 1$.
We fix all columns of $\mathbf A$ except for the $j$-th column and then average over the remaining one.
Since $z_j=\pm 1$ is invertible modulo $q$, the previous claim follows.

We note that
$
\sum_{\mathbf y \in \mathbb Z_q^n}
|S_{\mathbf A}(\mathbf y)|^2
=
\left|
\left\{
(\mathbf x,\mathbf x')
:
\mathbf A\mathbf x = \mathbf A\mathbf x'
\right\}
\right|.
$
Hence, taking the expectation over $\mathbf A$ gives:
$$
\begin{aligned}
\mathbb E_{\mathbf A}
\left[
\sum_{\mathbf y \in \mathbb Z_q^n}
|S_{\mathbf A}(\mathbf y)|^2
\right]
&=
2^m
+
\Pr_{\mathbf A}
\left[
\mathbf A\mathbf x=\mathbf A\mathbf x'
:
\mathbf x,\mathbf x' \in \{0,1\}^m,\ \mathbf x \neq \mathbf x'
\right]
\cdot
(2^{2m}-2^m)
\\
&=
2^m
+
\dfrac{2^{2m}-2^m}{q^n}.
\end{aligned}
$$

From this, we have
$$
\begin{aligned}
\Delta
&\leq
\sqrt{
q^n
\left(
\dfrac{1}{2^{2m}}
\left(
2^m+\dfrac{2^{2m}-2^m}{q^n}
\right)
-
\dfrac{1}{q^n}
\right)
}
\\
&=
\sqrt{
q^n
\left(
\dfrac{1}{2^m}
-
\dfrac{1}{2^m q^n}
\right)
}
\\
&\leq
\sqrt{\dfrac{q^n}{2^m}}
\\
&\leq
\dfrac{1}{q^{n/2}}
\\
&\leq
2^{-n},
\end{aligned}
$$
where the last inequality follows from $q>4$. This completes the proof.
\end{proof}

\begin{lemma}[Regularity Lemma~\cite{C:GenPeiVai08}]
\label{lem:4}
Let $n$ be a positive integer, let $q$ be a prime, and let
$m \ge 2n\log q$. Then, for all but a negligible fraction of
$\mathbf{A} \in \mathbb{Z}_q^{n\times m}$, and for any
$\sigma \ge \omega(\sqrt{\log m})$, we have 
$$\{\mathbf{Ae} \mod q:e \leftarrow \mathcal{D}_{\mathbb{Z}^m,\sigma}\} \approx_s \mathbf{U}_\mathbb{Z_q^n}$$
\end{lemma}


\begin{lemma}[From~\cite{EC:CasHofKilPei10}, adapted from Lemma~16~\cite{EPRINT:BonLauTas25}]
\label{lem:tail_bound}
Let $q \geq 2$ and $\mathbf{A} \in \mathbb{Z}_q^{n \times m}$ be a matrix
with $m > n$. Then, for $\mathbf{u} \in \mathbb{Z}_q^n$, it holds that
\[
    \Pr\left[
        \|\mathbf{x}\| > \sqrt{m}\sigma
        :
        \mathbf{x} \leftarrow \mathbf{A}_{\sigma}^{-1}(\mathbf{u})
    \right]
    \leq \operatorname{negl}(n).
\]
\end{lemma}

\begin{lemma}[From~\cite{C:MicPei12}]
\label{lem:gaget_matrix_sam}
For any integers $q \geq 2$, $n \geq 1$, $k = (\lceil \log q \rceil + 1)$ and $m = nk$, there is a primitive matrix
$\mathbf{G} \in \mathbb{Z}_q^{n \times m}$ such that the lattice $\Lambda^\perp(\mathbf{G})$ has a known basis $\mathbf{S} \in \mathbb{Z}^{m \times m}$ with $\|\widetilde{\mathbf{S}}\| \leq \sqrt{5}$ and $\|\mathbf{S}\| \leq \max\{\sqrt{5}, \sqrt{k}\}$.
\end{lemma}

\begin{lemma}[From~\cite{C:GenPeiVai08}]
\label{lem:uniform_mul_gaussian}
Let $n$ and $q$ be positive integers, where $q$ is prime, and let
$m \geq 2n\log q$. Then, except for at most a $2q^{-n}$ fraction of
matrices $\mathbf{A} \in \mathbb{Z}_q^{n \times m}$, for every
$s \geq \omega(\sqrt{\log m})$, the syndrome
$
\mathbf{u} = \mathbf{A}\mathbf{e} \bmod q
$
is statistically close to the uniform distribution over $\mathbb{Z}_q^n$,
where $\mathbf{e} \leftarrow D_{\mathbb{Z}^m,s}$.
\end{lemma}

\begin{lemma}[From~\cite{C:GenPeiVai08}]
\label{lem:5}
Let $n$ be a positive integer, let $q$ be a prime, and let
$m' = n(\lceil \log q \rceil + 1)$. For any
$\sigma \geq \omega(\sqrt{\log m'})$, if
$\mathbf{v} \leftarrow \mathbb{Z}_q^n$ and
$\mathbf{u} \leftarrow \mathbf{G}_{\sigma}^{-1}(\mathbf{v})$, then
$\mathbf{u}$ is statistically close to
$D_{\mathbb{Z}^{m'},\sigma}$.
\end{lemma}
\begin{proof}
By lemma~\ref{lem:1}, the smoothing pararamter $\eta_{\varepsilon}(\Lambda^\perp(\mathbf{G}) \leq \sqrt{5}\cdot \omega(\sqrt{\log n})$, which immediately implies that $\sigma \geq \eta_{\varepsilon}(\Lambda^\perp(\mathbf{G}))$. Finally, note that the columns of $\mathbf{G}$ generate $\mathbb{Z}_q^n$, by Lemma~\ref{lem:gaget_matrix_sam}, we derive the desirable result.
\end{proof}
\section{Hard Assumptions}

\begin{definition}[Learning with Errors -- LWE]

Ever since the first formal reduction from worst case to average case by Regev~\cite{J:Regev09}, Learning with Errors (LWE) has seen an intensive adoption and become a crucial building block of several cryptographic primitives.
At its core, LWE can be viewed as a \emph{noisy} linear algebra problem. In standard (noiseless) linear algebra over $\mathbb{Z}_q$, given a matrix $\mathbf{A} \in \mathbb{Z}_q^{m \times n}$ and a vector $\mathbf{b} = \mathbf{A}\mathbf{s}$, one can recover the secret $\mathbf{s} \in \mathbb{Z}_q^n$ by Gaussian elimination. LWE makes this computationally intractable by adding a small random error: the adversary sees $\mathbf{b} = \mathbf{A}\mathbf{s} + \mathbf{e}$, where $\mathbf{e}$ is a noise vector drawn from an error distribution $\chi$ (typically a discrete Gaussian $\mathcal{D}_{\mathbb{Z},\sigma}$ with $\sigma \ll q$). The noise is large enough to defeat algebraic attacks yet small enough relative to $q$ that it does not destroy the useful structure needed for cryptographic constructions.

Two hardness variants are studied. The \emph{search} variant asks to recover $\mathbf{s}$ given arbitrarily many samples $(\mathbf{a}_i, \mathbf{a}_i^\top\mathbf{s} + e_i)$. The \emph{decisional} variant, which suffices for the constructions in this thesis, asks only to \emph{distinguish} such samples from uniformly random pairs. Regev~\cite{J:Regev09} proved a (quantum) reduction from worst-case lattice problems (GapSVP and SIVP) to average-case decisional LWE, grounding its hardness in a well-studied geometric problem.

Formally, given a prime $q$, a positive integer $n$, and a distribution $\chi$ over $\mathbb{Z}_q$, the
$\mathrm{LWE}_{\mathbb{Z}_q,n,\chi}$ problem instance is associated with a challenge oracle $\mathcal{O}$ that behaves as follows:
\begin{itemize}
    \item $\mathcal{O}_{\mathbf{s}}$ outputs samples of the form
    $(\mathbf{u}_i,v_i):=(\mathbf{u}_i,\mathbf{u}_i^t \mathbf{s}+x_i)\in \mathbb{Z}_q^n\times \mathbb{Z}_q,$
    where $\mathbf{u}_i\leftarrow \mathbb{Z}_q^n$, $\mathbf{s}\leftarrow \mathbb{Z}_q^n$ is a uniformly random secret, and
    $x_i\in \mathbb{Z}_q$ is a fresh sample from $\chi$.

    \item $\mathcal{O}_{\Phi}$ outputs uniformly random samples from $\mathbb{Z}_q^n\times \mathbb{Z}_q$.
\end{itemize}

An algorithm $\mathcal{A}$ decides the $(\mathbb{Z}_q,n,\chi)$-LWE problem if, for some
$\mathbf{s}\leftarrow \mathbb{Z}_q^n$, the advantage of $\mathcal{A}$, defined as
\[
\mathrm{Adv}_{\mathrm{LWE}}[\mathcal{A}]
=
\left|
\Pr[\mathcal{A}^{\mathcal{O}_{\mathbf{s}}}=1]
-
\Pr[\mathcal{A}^{\mathcal{O}_{\Phi}}=1]
\right|,
\]
is non-negligible in $\lambda$.
\end{definition}

\begin{definition}[$\kappa$-Uniform Distribution~\cite{AC:ALLW25}]
Let $\varphi,\kappa,m,\log q \in \mathrm{poly}(\lambda)$ with $\kappa \le m$.
For any $\mathbf{U} \in \mathbb{Z}_q^{m\times \kappa}$, denote
$\mathbf{U}^{\perp} = \left\{\mathbf{x} \in \mathbb{Z}_q^m : \mathbf{x}^t \cdot \mathbf{U} = \mathbf{0}^t \bmod q \right\}.$
The $\kappa$-uniform distribution with parameters $\mathbf{U},\mathbf{R},q,\chi$ is
\[
\mathcal{U}(\mathbf{U}^{\perp}) + \mathcal{D}_{\mathbb{Z},\chi}^{m}
:=
\left\{
\mathbf{b}
\;\middle|\;
\mathbf{x}\leftarrow \mathbf{U}^{\perp};
\;
\mathbf{e}\leftarrow \mathcal{D}_{\mathbb{Z},\chi}^{m};
\;
\mathbf{b}:=\mathbf{x}+\mathbf{e}
\right\}.
\]
\end{definition}

\begin{definition}[$\kappa$-LWE Assumption With Varied Hints~\cite{AC:ALLW25}]
Let $n,m,\log q,\kappa,\chi,\{s_{\max}(\mathbf{S}_i)\}_i \in \mathrm{poly}(\lambda)$ with $n,\kappa \le m$.
Let $\{\mathbf{S}_i\}_i$ be a set of arbitrary full-rank matrices
$\mathbf{S}_i \in \mathbb{Z}^{ m \times  m}$.
The $\kappa$-LWE$_{\mathbb{Z}_q,n,m,\{\mathbf{S}_i\}_i,\chi}$ assumption states that for any PPT $\mathcal{A}$,
the following distributions over $(\mathbf{A},\mathbf{U},\mathbf{b})$:
\[
\left\{
\begin{aligned}
\mathbf{A} &\leftarrow \mathbb{Z}_q^{n\times m} \\
\mathbf{U} &\leftarrow \prod_{i=0}^{\kappa-1} \mathcal{D}_{\Lambda_q^{\perp}(\mathbf{A}),\mathbf{S}_i} \\
\mathbf{s} &\leftarrow \mathbb{Z}_q^{n},\;\mathbf{e}\leftarrow \mathcal{D}_{\mathbb{Z},\chi}^{m} \\
\mathbf{b}^{t} &:= \mathbf{s}^{t}\cdot \mathbf{A} + \mathbf{e}^{t} \bmod q
\end{aligned}
\right\}
\approx_c
\left\{
\begin{aligned}
\mathbf{A} &\leftarrow \mathbb{Z}_q^{n\times m} \\
\mathbf{U} &\leftarrow \prod_{i=0}^{\kappa-1} \mathcal{D}_{\Lambda_q^{\perp}(\mathbf{A}),\mathbf{S}_i} \\
\mathbf{b} &\leftarrow \mathcal{U}(\mathbf{U}^{\perp}) + \mathcal{D}_{\mathbb{Z},\chi}^{m}
\end{aligned}
\right\}.
\]
\end{definition}

\begin{theorem}[Adapted from~\cite{AC:ALLW25}]
\label{theo:k_LWE}
Let $n, m, \kappa, q \in \mathbb{Z}_{\geq 1}$, let $\beta, \beta' > 0$ be reals,
let $\chi, \chi'$ be noise distributions, and let
$\{\sigma_i, \sigma_i'\}_{i \in [\kappa]}$ be positive reals satisfying the
following conditions:
\begin{itemize}
  \item $\kappa \geq 100$, \quad and \quad $\sigma_i > 0,\ \sigma_i' > 0$
        for all $i \in [\kappa]$;
  \item $m \geq \max\!\Big\{\, \Omega\!\big(\kappa \log(\kappa \max_i \sigma_i)\big),\
        \Omega(n \log q) \,\Big\}$;
  \item $q > \max_i(\sigma_i') \cdot \sqrt{\log m}$ \quad and \quad $q$ is prime;
  \item $\displaystyle \min_i(\sigma_i)
        > \max\!\Big\{\, \Omega\!\big(\sqrt{\kappa m \log m}\big),\
        \max_i(\sigma_i')^{\,\kappa/(m+\kappa)} \,\Big\}$;
  \item $\displaystyle \min_i(\sigma_i')
        \geq \Omega\!\Big( \kappa \sqrt{m}\, \max_i(\sigma_i)^2\,
        \log^{3/2}\!\big(\kappa m \max_i(\sigma_i)\big) \Big)$.
\end{itemize}

For each $i \in [\kappa]$, define the block-diagonal matrix
$
  \mathbf{S}_i \;=\; \operatorname{diag}\!\big(\sigma_i \mathbf{I},\, \sigma_i' \mathbf{I}\big).
$
If, in addition,
$
  \beta \;>\; \Omega\!\big( m^{3/2} \max_i(\sigma_i') \cdot \beta' \big)$
  and
$\chi' \;>\; \Omega\!\big( m^{3/2} \max_i(\sigma_i') \cdot \chi \big),
$
then there exist a \textup{PPT} reduction from
        $\mathrm{LWE}_{\mathbb{Z}_q,\, n,\, m,\, \chi}$
        to
        $\kappa\text{-}\mathrm{LWE}_{\mathbb{Z}_q,\, n,\, m+\kappa,\, \{\mathbf{S}_i\}_i,\, \chi'}$.

\end{theorem}
To answer the ciphertext challenge, the simulator relies on the following lemma to simulate the challenge ciphertext. This follows from the standard noise flooding/noise drowning lemma
\cite{GKPV10}, which states that a sufficiently wide discrete Gaussian
statistically hides any bounded additive shift.

\begin{lemma}[Noise flooding Lemma~\cite{GKPV10}]
\label{lem:̉noise_flooding}
Let $\lambda$ be a security parameter. Take any $e \in \mathbb{Z}$ where
$|e| \leq B$. Suppose
$
\chi \geq B \cdot \lambda^{\omega(1)}.
$
Then
\[
\triangle\left[\left\{ z : z \leftarrow D_{\mathbb{Z},\chi} \right\},
\left\{ z + e : z \leftarrow D_{\mathbb{Z},\chi} \right\}\right]
\leq negl(\lambda)\]
\end{lemma}

\noindent\section{Gadget trapdoors}

A trapdoor function is a fundamental component of public-key cryptosystems. For cryptographic applications, it must satisfy three main properties. First, the trapdoor should be sampled from a prescribed distribution. Second, there should exist an efficient preimage-sampling algorithm that uses the trapdoor as extra knowledge to recover a preimage of the one-way function without leaking any information about the secret. Finally, the output distribution of the aforementioned algorithm must also follow a prescribed distribution. 

\begin{definition} (Adapted from \cite{C:MicPei12}) For positive integers $n,q \in \mathbb{N}$, let 
$\mathbf{G}_n = \mathbf{I}_n \otimes \mathbf{g}^{\top} \in \mathbb{Z}_q^{n \times m'}$ 
be the gadget matrix, where 
$\mathbf{g}^{\top} = [1,2,\ldots,2^{\lceil \log q \rceil}]$ 
and $m' = n(\lceil \log q \rceil + 1)$.
\end{definition}


\begin{lemma}[Adapted from \cite{C:MicPei12}]
\label{lem:gadget_trapdoor}
Let $n,m,q$ be lattice parameters and
$m' = n(\lceil \log q \rceil + 1).$
Then, there exist efficient algorithms $(\mathsf{TrapGen}, \mathsf{SamplePre})$ such that:
\begin{itemize}
    \item $\mathsf{TrapGen}(1^n, q, m) \to (\mathbf{A}, \mathbf{T}_{\mathbf{A}})$:
    On input the lattice dimension $n$, the modulus $q$, and the number of samples $m$, the trapdoor-generation algorithm outputs a matrix
    $\mathbf{A} \in \mathbb{Z}_q^{n \times m}$
    together with a gadget trapdoor
    $\mathbf{T}_{\mathbf{A}} \in \mathbb{Z}_q^{m \times m'}$

    \item $\mathsf{SamplePre}(\mathbf{A}, \mathbf{T}_{\mathbf{A}}, \mathbf{v}, \sigma) \to \mathbf{u}$:
    On input a matrix
    $\mathbf{A} \in \mathbb{Z}_q^{n \times m},$
    a trapdoor
    $\mathbf{T}_{\mathbf{A}} \in \mathbb{Z}_q^{m \times m'},$
    a target vector
    $\mathbf{v} \in \mathbb{Z}_q^n,$
    and a Gaussian width parameter $\sigma$, the preimage-sampling algorithm outputs a vector
    $\mathbf{u} \in \mathbb{Z}_q^m.$
\end{itemize}

For all $m \ge O(n \log q)$, the above algorithms satisfy the following properties:
\begin{itemize}
    \item \textbf{Trapdoor distribution:}
    The matrix $\mathbf{A}$ output by $\mathsf{TrapGen}(1^n,q,m)$ is statistically close to uniform. Specifically, if
    $(\mathbf{A},\mathbf{T}_{\mathbf{A}}) \leftarrow \mathsf{TrapGen}(1^n,q,m)$ and
    $\mathbf{A}' \leftarrow \mathbb{Z}_q^{n\times m}$, then we have
    $\Delta[\mathcal{D}(\mathbf{A}),\mathcal{D}(\mathbf{A}')] \le 2^{-n}$.
    Moreover,
    $\mathbf{A}\mathbf{T}_{\mathbf{A}} = \mathbf{G}$
    and
    $\|\mathbf{T}_{\mathbf{A}}\|_{\infty} = 1$.

    \item \textbf{Preimage distribution.}
        Suppose $\mathbf{T}_{\mathbf{A}}$ is a gadget trapdoor for
        $\mathbf{A} \in \mathbb{Z}_q^{n \times m}$, i.e.,
        $\mathbf{A}\mathbf{T}_{\mathbf{A}} \equiv \mathbf{G}_n \pmod q$.
        For every
        $
        \sigma \ge
        \sqrt{mm'}\,\|\mathbf{T}_{\mathbf{A}}\|_{\infty}
        \cdot \omega(\sqrt{\log n})
        $
        and every target $\mathbf{v} \in \mathbb{Z}_q^n$, we have
        $
        \Delta\!\left[
        \mathsf{SamplePre}
        (\mathbf{A},\mathbf{T}_{\mathbf{A}},\mathbf{v},\sigma),
        \mathcal{D}_{\Lambda_q^{\mathbf{v}}(\mathbf{A}),\sigma}
        \right]
        \le 2^{-n}.
        $
\end{itemize}

For a matrix
$\mathbf{V} := [\mathbf{v}_1,\ldots,\mathbf{v}_{\ell}] \in \mathbb{Z}_q^{n\times \ell}$,
we define $\mathsf{SamplePre}(\mathbf{A},\mathbf{T}_{\mathbf{A}},\mathbf{V},\sigma)$ as the algorithm that outputs the matrix whose $i$-th column is
$\mathsf{SamplePre}(\mathbf{A},\mathbf{T}_{\mathbf{A}},\mathbf{v}_i,\sigma)$.
\end{lemma}



\noindent\section{Ajtai trapdoors} 

It is first introduced in~\cite{EC:ABB10} and has subsequently been expanded and used in several works.


\begin{definition}[Ajtai Trapdoor~\cite{C:Ajt96}]
Let $n,m,q$ be lattice parameters and let $\mathbf{A} \in \mathbb{Z}_q^{n \times m}$ be a matrix. 
A matrix $\mathbf{R} \in \mathbb{Z}^{m \times m}$ is an \emph{Ajtai trapdoor} for $\mathbf{A}$ if 
$\mathbf{A}\mathbf{R} = \mathbf{0} \bmod q$ and $\mathbf{R}$ is full rank over $\mathbb{R}$. 
\end{definition}


\begin{definition}


Let $\mathbf A \in \mathbb Z_q^{n \times m}$, 
$\mathbf M_1 \in \mathbb Z_q^{n \times m_1}$, and let $\sigma$ be a Gaussian parameter. 
For any target vector $\mathbf u \in \mathbb Z_q^n$, the efficient probabilistic algorithm
$
    \mathsf{SampleLeft}(\mathbf A,\mathbf M_1,\mathbf T_{\mathbf A},\mathbf u,\sigma)
$
of~\cite{EC:ABB10} uses a short basis $\mathbf T_{\mathbf A}$ of 
$\Lambda_q^\perp(\mathbf A)$ to sample a vector
$
    \mathbf v \in \Lambda_q^{\mathbf u}(\mathbf F),
    \text{where } \mathbf F = [\mathbf A \mid \mathbf M_1],
$
from a distribution statistically close to 
$D_{\Lambda_q^{\mathbf u}(\mathbf F),\sigma}$.


\end{definition}

\begin{lemma}[SampleLeft, Lemma 17 of~\cite{EC:ABB10}]
\label{lem:sampleleft}
Let $q>2$, $m>n$, and
$
\sigma > \|\widetilde{\mathbf T}_{\mathbf A}\|_{\mathsf{GS}}
\cdot \omega(\sqrt{\log(m+m_1)}).
$
Then $\mathsf{SampleLeft}(\mathbf A,\mathbf M_1,\mathbf T_{\mathbf A},\mathbf u,\sigma)$ outputs a vector distributed statistically close to $\mathbf F_\sigma^{-1}(\mathbf u)$, where $\mathbf F=[\mathbf A\mid \mathbf M_1]$.
\end{lemma}

The input trapdoor for this sampling algorithm is not limited only to Ajtai trapdoor but also to the gadget trapdoor, which was introduced in the concurrent work~\cite{C:MicPei12}. This interesting property is discussed further in Appendix~\ref{ap:3}


\section{Explainable SampleLeft}
In cryptographic construction, a random sampler is essential to produce outputs from a desirable distribution. However, statistical closeness is barely sufficient since sampling functions also need to be programmable for security proof. To address this issue, a new notion called 'explainability' was proposed in \cite{TCC:LuWat22}, which particularly allows reversing sample an $r$ given input $x$ and $\mathsf{Sample}(r) = x$
\begin{definition}[Explainable $\mathsf{SampleLeft}$]
Let $n,m,m_1,q$ be lattice parameters. A $(\rho,\sigma_{\mathsf{loss}})$-explainable $\mathsf{SampleLeft}$ sampler is a pair of efficient algorithms
$
\Pi_{\mathsf{SL}}=(\mathsf{SampleLeft},\mathsf{ExplainSL})
$
with the following syntax:
\begin{itemize}
    \item $\mathsf{SampleLeft}(\mathbf A,\mathbf M_1,\mathbf T_{\mathbf A},\mathbf u,\sigma;\mathsf{rb}) \to \mathbf v$:
    on input $\mathbf A\in\mathbb Z_q^{n\times m}$, $\mathbf M_1\in\mathbb Z_q^{n\times m_1}$, a trapdoor $\mathbf T_{\mathbf A}$  for $\mathbf A$, $\mathbf u\in\mathbb Z_q^n$, Gaussian width $\sigma$, and randomness $\mathsf{rb}\in\{0,1\}^{\rho}$, outputs $\mathbf v\in\mathbb Z_q^{m+m_1}$.

    \item $\mathsf{ExplainSL}(1^\kappa,\mathbf A,\mathbf M_1,\mathbf T_{\mathbf A},\mathbf u,\mathbf v,\sigma) \to \mathsf{rb}$:
    on input the same public parameters and a vector $\mathbf v$, outputs randomness $\mathsf{rb}\in\{0,1\}^{\rho}$.
\end{itemize}

Let $\mathbf F=[\mathbf A\mid \mathbf M_1]$. Except with negligible probability, for all admissible parameters, matrices, trapdoors, targets $\mathbf u$, and widths $\sigma$ satisfying
$
\|\mathbf T_{\mathbf A}\|_{\infty}\cdot \sigma_{\mathsf{loss}}(\lambda,n,m,m_1,q)
\leq
\sigma
\leq
2^\lambda,
$
the following hold.

\textbf{Correctness.} The distributions
$$
\left\{
\mathbf v \leftarrow
\mathsf{SampleLeft}(\mathbf A,\mathbf M_1,\mathbf T_{\mathbf A},\mathbf u,\sigma;\mathsf{rb})
:
\mathsf{rb}\leftarrow\{0,1\}^{\rho}
\right\}
\quad\text{and}\quad
\left\{
\mathbf v \leftarrow \mathbf F_{\sigma}^{-1}(\mathbf u)
\right\}
$$
are statistically close. Moreover, every output $\mathbf v$ satisfies $\mathbf F\mathbf v=\mathbf u$.

\textbf{Explainability.} For every precision parameter $\kappa\in\mathbb N$, the statistical distance between the following two experiments is at most $1/\kappa+\mathsf{negl}(\lambda)$:
$$
\begin{array}{ll}
\mathsf{DSampleLeft}:&
\mathsf{rb}\leftarrow\{0,1\}^{\rho},\
\mathbf v\leftarrow
\mathsf{SampleLeft}(\mathbf A,\mathbf M_1,\mathbf T_{\mathbf A},\mathbf u,\sigma;\mathsf{rb});
\ \text{output }(\mathbf v,\mathsf{rb}).
\\[0.5em]
\mathsf{DExplainLeft}:&
\mathsf{rb}'\leftarrow\{0,1\}^{\rho},\
\mathbf v\leftarrow
\mathsf{SampleLeft}(\mathbf A,\mathbf M_1,\mathbf T_{\mathbf A},\mathbf u,\sigma;\mathsf{rb}');
\\
&
\mathsf{rb}\leftarrow
\mathsf{ExplainSL}(1^\kappa,\mathbf A,\mathbf M_1,\mathbf T_{\mathbf A},\mathbf u,\mathbf v,\sigma);
\ \text{output }(\mathbf v,\mathsf{rb}).
\end{array}
$$
\end{definition}

\begin{theorem}
\label{the:explainable}
The $(\rho,\sigma_{\mathsf{loss}})$-explainable described in construction~\ref{ap:2}
$\mathsf{SampleLeft}$ sampler $\Pi_{\mathsf{SL}}$ for
$
    \rho = \operatorname{poly}(\lambda,n,m,m_1,\log(q))
$
and
$
    \sigma_{\mathsf{loss}}(\lambda,n,m,m_1,q)
    =
    18(m+m_1)^{3/2}
    \log((m+m_1)\lambda)
    \log(\log(q)).
$
\end{theorem}



\section{Shifted multi-preimage trapdoor sampler}

My protocol employs  the multi-preimage protocol proposed by Waters,Wee,and Wu ~\cite{EC:WatWeeWu25}, which is then modified to be suitable for the construction in ~\cite{EPRINT:BonLauTas25}, Therefore, we 
\label{def:shifted-multi-preimage-sampler}
Let $\ell$ be the batch size and $d,n,t,q,\sigma$ be functions of
$\lambda$ and $\ell$. An
$(n,t,q,\sigma)$-\emph{shifted multi-preimage trapdoor sampler} consists of
efficient algorithms
$
    \mathsf{Gen},
    \mathsf{Expand},
    \mathsf{SampleMultPre},
$
defined as follows.
\begin{itemize}
    \item 
    $\mathsf{Gen}(1^\lambda,1^\ell,1^d)\to \mathsf{crs}$ is a randomized
    generation algorithm that outputs a common reference string $\mathsf{crs}$.

    \item
    $\mathsf{Expand}(1^\lambda,1^\ell,1^d,\mathsf{crs},(u_i)_{i\in[\ell]})
    \to (\mathbf A_1,\ldots,\mathbf A_\ell,\mathsf{td})$
    is a deterministic expansion algorithm. It takes $\ell$ distinct bit
    vectors $(u_i)_{i\in[\ell]}$ and outputs matrices
    $
        \mathbf A_1,\ldots,\mathbf A_\ell \in \mathbb Z_q^{n\times t}
    $
    together with a matrix-trapdoor tuple $\mathsf{td}$.

    \item
    $\mathsf{SampleMultPre}(\mathsf{td},\mathbf t_1,\ldots,\mathbf t_\ell)
    \to (\boldsymbol\pi_1,\ldots,\boldsymbol\pi_\ell,\mathbf c)$
    is a randomized sampling algorithm. Given a matrix-trapdoor tuple
    $\mathsf{td}$ and target vectors
    $\mathbf t_1,\ldots,\mathbf t_\ell\in\mathbb Z_q^n$, it outputs a shift
    $
        \mathbf c\in\mathbb Z_q^n
    $
    and preimages.
    $
        \boldsymbol\pi_1,\ldots,\boldsymbol\pi_\ell\in\mathbb Z^t .
    $
\end{itemize}
A shifted multi-preimage trapdoor sampler has the following properties:

\begin{itemize}
    \item \textbf{Correctness.} For all $\lambda,\ell,d\in\mathbb N$, all
$\mathsf{crs}$ in the support of $\mathsf{Gen}(1^\lambda,1^\ell,1^d )$,
and all target vectors
$\mathbf t_1,\ldots,\mathbf t_\ell\in\mathbb Z_q^n$, let
$
    (\mathbf A_1,\ldots,\mathbf A_\ell,\mathsf{td})
    :=
    \mathsf{Expand}
    (1^\lambda,1^\ell,1^d,\mathsf{crs},(u_i)_{i\in[\ell]}).
$
Then
\[
\Pr\left[
    \mathbf A_i\boldsymbol\pi_i=\mathbf t_i+\mathbf c
    \ \forall i\in[\ell]
    \;:\;
    (\boldsymbol\pi_1,\ldots,\boldsymbol\pi_\ell,\mathbf c)
    \leftarrow
    \mathsf{SampleMultPre}
    (\mathsf{td},\mathbf t_1,\ldots,\mathbf t_\ell)
\right]
=1 .
\]
\item \textbf{Preimage distribution.}
For every polynomial $\ell=\ell(\lambda)$, except with negligible probability
over
$
    \mathsf{crs}\leftarrow \mathsf{Gen}(1^\lambda,1^\ell,1^d),
$
and
$
    (\mathbf A_1,\ldots,\mathbf A_\ell,\mathsf{td})
    :=
    \mathsf{Expand}
    (1^\lambda,1^\ell,1^d,\mathsf{crs},(u_i)_{i\in[\ell]}),
$
the statistical distance between the following two distributions is
$\mathsf{negl}(\lambda)$:
\[
\left\{
    (\boldsymbol\pi_1,\ldots,\boldsymbol\pi_\ell,\mathbf c)
    \leftarrow
    \mathsf{SampleMultPre}
    (\mathsf{td},\mathbf t_1,\ldots,\mathbf t_\ell)
\right\}
\]
and
\[
\left\{
    (\boldsymbol\pi_1,\ldots,\boldsymbol\pi_\ell,\mathbf c)
    \ \middle|\
    \mathbf c\leftarrow \mathbb Z_q^n,\;
    \boldsymbol\pi_i
    \leftarrow
    (\mathbf A_i^{-1})_\sigma(\mathbf t_i+\mathbf c)
    \ \forall i\in[\ell]
\right\}.
\]

\item \textbf{Somewhere programmability.}
There exists an efficient algorithm $\mathsf{GenProg}$ such that, for every
polynomial $\ell=\ell(\lambda)$ and every collection of $\ell$ distinct bit
vectors $(u_i)_{i\in[\ell]}$, the following hold.
\begin{itemize}
     

\item For every $i\in[\ell]$ and every
$\widetilde{\mathbf A}_i\in\mathbb Z_q^{n\times t}$,
\[
\Pr\left[
    \mathbf A_i=\widetilde{\mathbf A}_i
    \ \middle|\
    \begin{array}{l}
    \widetilde{\mathsf{crs}}
    :=
    \mathsf{GenProg}
    (1^\lambda,1^\ell,1^d,u_i,\widetilde{\mathbf A}_i),\\[0.3em]
    (\mathbf A_1,\ldots,\mathbf A_\ell,\mathsf{td})
    :=
    \mathsf{Expand}
    (1^\lambda,1^\ell,1^d,\widetilde{\mathsf{crs}},
    (u_i)_{i\in[\ell]})
    \end{array}
\right]
=1 .
\]

\item The statistical distance between
$
    \left\{
        \mathsf{crs}\leftarrow \mathsf{Gen}(1^\lambda,1^\ell)
    \right\}
$
and
\[
    \left\{
        \widetilde{\mathsf{crs}}
        \ \middle|\
        \mathbf A_i\leftarrow \mathbb Z_q^{n\times t},\;
        \widetilde{\mathsf{crs}}
        :=
        \mathsf{GenProg}
        (1^\lambda,1^\ell,1^d,u_i,\mathbf A_i)
    \right\}
\]
is $\mathsf{negl}(\lambda)$.
\end{itemize}
\item\textbf{Local expansion.}
There exists an efficient deterministic algorithm $\mathsf{ExpandLocal}$ such
that
$
    \mathsf{ExpandLocal}(1^\lambda,\mathsf{crs},u_i)
    \to \mathbf A_i\in\mathbb Z_q^{n\times t}.
$
Moreover, for all $\lambda,\ell\in\mathbb N$ and all $\mathsf{crs}$, if
\[
    (\mathbf A_1,\ldots,\mathbf A_\ell,\mathsf{td})
    :=
    \mathsf{Expand}
    (1^\lambda,1^\ell,1^d,\mathsf{crs},(u_i)_{i\in[\ell]}),
\]
then for every $i\in[\ell]$,
$
    \mathsf{ExpandLocal}(1^\lambda,\mathsf{crs},u_i)
    =
    \mathbf A_i .
$\\
Finally, $\mathsf{ExpandLocal}(1^\lambda,\mathsf{crs},u_i)$ reads only
$\mathsf{poly}(\lambda,\log \ell)$ bits of $\mathsf{crs}$ and runs in time
$\mathsf{poly}(\lambda,\log \ell)$.

\end{itemize}
\begin{lemma}[Adapted from lemma 2 of~\cite{EPRINT:BonLauTas25}]
\label{lem:smps}
Let $n \geq \lambda$, $d=d(\lambda)$, $n=n(\lambda,\ell,d)$,
$q=q(\lambda,\ell,d)$, $m=3n\lceil \log q\rceil$, and
$t=\widetilde{m}(d+1)$. For any
$
    \sigma \geq (\ell t+\widetilde{m})\log(\ell n),
$
the construction in~\ref{ap:1} instantiated with any distinct labels
$u_1,\ldots,u_\ell \in \{0,1\}^d$, yields an
$(n,t,q,\sigma)$-shifted multi-preimage trapdoor sampler.
Moreover, this sampler satisfies perfect somewhere programmability,
transparent setup, local expansion, and simulatable openings. Its CRS
has size
$
    nt\log(q)=n\widetilde{m}(d+1)\log(q),
$
and the algorithm $\mathsf{Expand}$ runs in time
$\operatorname{poly}(n,m,\ell,d,\log q)$.
\end{lemma}


\section{Partial Lattice Trapdoors}
\begin{definition} (Partial Lattice Trapdoors) A $(k,N)$-threshold partial lattice trapdoor scheme over $\mathbb{Z}_q$
is an extension of a (full) lattice trapdoor scheme
$(\mathsf{TrapGen}, \mathsf{SampPre})$ over $\mathbb{Z}_q$ with the PPT algorithms
$(\mathsf{PTrapGen}, \mathsf{PSampPre}, \mathsf{Rec})$:

$
 \mathsf{PTrapGen}(\mathbf{A}, \mathsf{td}, j) \rightarrow \mathsf{ptd}_j:
$
The partial trapdoor generation algorithm inputs a matrix $\mathbf{A}$,
its (full) trapdoor $\mathsf{td}$ and an index $j \in [N]$, and generates
a partial trapdoor $\mathsf{ptd}_j$ of $\mathbf{A}$ for $j$.

$
 \mathsf{PSampPre}(\mathbf{A}, \mathsf{ptd}_j, T, \mathbf{y}, \sigma) \rightarrow \mathbf x_j:
$
The partial preimage sampling algorithm, given a partial trapdoor
$\mathsf{ptd}_j$ for index $j$, a set $T \subset [N]$, a target image $\mathbf{y}$,
and a Gaussian parameter $\sigma$, samples a partial preimage $\mathbf x_j$ for $j$.

$
\mathsf{Rec}\bigl((\mathbf x_j)_{j \in T}\bigr) \rightarrow \mathbf x:
$
The reconstruction algorithm, given a tuple of partial preimages
$(\mathbf x_j)_{j \in T}$, reconstructs a preimage $\mathbf x$.

\vspace{1em}

\begin{itemize}
    \item  \textbf{Correctness}:
Let $N,n,m,\widetilde{m},k \in \mathbb{N}$ with $k \leq N$ and $n \leq m$,
and $\widetilde{s}, \beta, \sigma > 0$.
A $(k,N)$-threshold partial lattice trapdoor scheme is said to be
$(n,m,\widetilde{s},\beta,\sigma)$-correct if the underlying lattice trapdoor
scheme $(\mathsf{TrapGen}, \mathsf{SampPre})$ is $(n,m,\widetilde{s})$-correct,
and for any $k$-subset $T \subseteq [N]$ and vector $\mathbf y \in \mathbb{Z}_q^n$, it holds that

\[
\Pr\!\left[
\begin{array}{l}
    \mathbf{A}\cdot \mathbf{x} \equiv \mathbf{y} \pmod q;\\
    \qquad \|\mathbf{x} \| \le \beta
\end{array}
\;\middle|\;
\begin{array}{l}
    (\mathbf{A},\mathsf{td})
        \leftarrow
        \mathsf{TrapGen}(1^\lambda,1^n,1^m), \\[0.3em]
    \mathsf{ptd}_j
        \leftarrow
        \mathsf{PTrapGen}(\mathbf{A},\mathsf{td},j)
        \quad \forall j \in [N], \\[0.3em]
    \mathbf{x} _j
        \leftarrow
        \mathsf{PSampPre}(\mathbf{A},\mathsf{ptd}_j,T,\mathbf{y} ,\sigma)
        \quad \forall j \in T, \\[0.3em]
    \mathbf{x} 
        \leftarrow
        \mathsf{Rec}\bigl((\mathbf{x} _j)_{j\in T}\bigr)
\end{array}
\right]
\ge
1-\mathsf{negl}(\lambda)
\]

\item \textbf{One-wayness and Indistinguishability}:
Let $N,n,m,k \in \mathbb{N}$ with $k \leq N$ and $n \leq m$, and
$\beta,\chi_0,\chi_1,\sigma > 0$. Denote
$
\mathsf{par}_{\mathsf{ow}} = (n,m,\beta,\sigma)
\text{ and } 
\mathsf{par}_{\mathsf{IND}} = (n,m,\chi_0,\chi_1,\sigma).
$
A $(k,N)$-threshold partial lattice trapdoor scheme is said to be $\mathsf{par}_{\mathsf{OW}}$-one-way and $\mathsf{par}_{\mathsf{IND}}$indistinguishable if, 
for any PPT adversary $\mathcal{A}$, it holds that
\[
\left\{
\begin{aligned}
\Pr\!\left[
\mathsf{Exp}^{\mathsf{OW}}_{\Pi,\mathcal{A},\mathsf{par}_{\mathsf{OW}}}(1^\lambda)
= 1
\right]
\le
\mathsf{negl}(\lambda),
\\[1.2em]
\left|
\Pr\!\left[
\mathsf{Exp}^{\mathsf{IND},0}_{\Pi,\mathcal{A},\mathsf{par}_{\mathsf{IND}}}(1^\lambda)
= 1
\right]
-
\Pr\!\left[
\mathsf{Exp}^{\mathsf{IND},1}_{\Pi,\mathcal{A},\mathsf{par}_{\mathsf{IND}}}(1^\lambda)
= 1
\right]
\right|
&\le
\mathsf{negl}(\lambda).
\end{aligned}
\right.
\]
\end{itemize}

\end{definition}

\begin{theorem}[Adapted from Theorem 2 of \cite{AC:ALLW25}]
Let parameters be as in Table~\ref{tab:parameters}, in particular
$
    \beta \ge k \cdot \sigma \cdot \sqrt{2mk\varphi}.
$
The partial lattice trapdoor in Construction~\ref{ap:4} is
$
    (nk,2mk,\widetilde{s},\beta,\sigma)$-correct.

\end{theorem}
\begin{theorem}[Adapted from Theorem 3 of \cite{AC:ALLW25}]
\label{theo:onewayness_partitial}
Let parameters be as in Table~\ref{tab:parameters}. The partial trapdoor scheme in Construction~\ref{ap:4} is
$(nk,2mk,\beta',\sigma)$-one-way if the
$\kappa$-$\mathsf{SIS}_{\mathbb{Z}_q,n,2mk,\{\widetilde{\mathbf{S}}_i\}_i,\alpha}$
assumption holds for $\kappa = 2m\cdot(k-1)$.
\end{theorem}

\begin{theorem}[Adapted from Theorem 4 of \cite{AC:ALLW25}]
\label{theo:indistinguishibity_oneway}
Let parameters be as in Table~\ref{tab:parameters}. The partial lattice trapdoor scheme in in Construction~\ref{ap:4} has
$(nk,2gmk,\chi_0,\chi_1,\sigma)$-indistinguishability if, for
$\kappa = 2m\cdot(k-1)$, the
$\kappa$-$\mathsf{LWE}_{\mathbb{Z}_q,n,2mk,\{\operatorname{diag}(\widetilde{\mathbf{S}}_i)\}_i,\chi_0}$
assumption holds.
\end{theorem}

Although the construction just need the indistinguishable property for proving formal security, I still include the one-wayness for completeness. 

% \end{document}